\chapter*{Abstract}\addcontentsline{toc}{chapter}{Abstract}
Mesospheric ozone and the 1.27~$\mu$m emission of \singlet\ are chemically connected, but a stationary interpretation can be inadequate near sunrise. This thesis develops a reproducible temporal forward model over 50--100~km, driven by UTC date/location solar geometry, spherical attenuation and either a frozen reference or prescribed dynamic MSIS atmosphere. Seven species evolve at 51 levels: O, O$_3$, H, OH, HO$_2$, H$_2$O$_2$ and $\Delta$, giving 357 ODEs. Only O($^1$D), $B_0$ and $B_1$ remain quasi-steady.

The original stationary radical partition lost a physical root during dawn evolution, and algebraic peroxide became singular in a dark limit. Their dynamic promotion preserves accepted reaction rates and event budgets while relaxing the incompatible stationarity assumptions. BDF integration is checked against tighter BDF and Radau; maximum relevant all-species differences are approximately $0.424\%$ and $0.00946\%$, respectively. Independent atmospheric/radiative interpolation refinements meet an approximately $0.5\%$ output target in the tested reference.

The final campaign describes reference dawn, four seasons and three latitudes under quiet activity and approximate noon initialization. Upper-level ozone decreases substantially while singlet oxygen grows, with a non-monotonic ozone response near 80~km. Seasonal ordering depends on height. An identical-initial-state frozen/dynamic comparison reaches maximum relevant ozone and singlet differences of $69.21\%$ and $24.61\%$. Atomic initial-condition variants also produce large upper-level sensitivities.

The result is a numerically validated finite-horizon mechanistic framework, not an observationally validated climatology. Approximate initialization, prescribed reservoirs/exterior ozone, quiet activity and absent transport limit atmospheric interpretation. Improved initial-state constraints, transport and compatible observational comparisons are prioritized future work.

\chapter*{Resumen}\addcontentsline{toc}{chapter}{Resumen}
El ozono mesosf\'erico y la emisi\'on de \singlet\ a 1,27~$\mu$m est\'an relacionados qu\'imicamente, pero una interpretaci\'on estacionaria puede resultar inadecuada cerca del amanecer. Este trabajo desarrolla un modelo temporal reproducible entre 50 y 100~km, con geometr\'ia solar basada en fecha UTC y localizaci\'on, atenuaci\'on esf\'erica y atm\'osfera de referencia congelada o background MSIS prescrito y variable. Se integran siete especies en 51 alturas: O, O$_3$, H, OH, HO$_2$, H$_2$O$_2$ y $\Delta$, con 357 ecuaciones diferenciales. Solo O($^1$D), $B_0$ y $B_1$ permanecen en aproximaci\'on cuasiestacionaria.

La partici\'on estacionaria original de radicales perdi\'o su ra\'iz f\'isica durante el amanecer, y la eliminaci\'on algebraica del per\'oxido se hizo singular en un l\'imite oscuro. Su promoci\'on a especies din\'amicas conserva las reacciones, constantes y presupuestos aceptados. La integraci\'on BDF se contrasta con BDF m\'as estricto y Radau; las diferencias m\'aximas relevantes entre las siete especies son aproximadamente 0,424\% y 0,00946\%. Los refinamientos independientes de la interpolaci\'on atmosf\'erica y radiativa cumplen el objetivo pr\'actico de aproximadamente 0,5\% en el caso certificado.

La campa\~na final estudia el amanecer de referencia, cuatro estaciones y tres latitudes con actividad quiet y bootstrap aproximado desde el mediod\'ia. El ozono disminuye fuertemente en alturas superiores mientras crece el ox\'igeno singlete; cerca de 80~km aparece una respuesta no mon\'otona. La ordenaci\'on estacional depende de la altura. Con id\'entico estado inicial, las atm\'osferas congelada y variable producen diferencias m\'aximas de 69,21\% en ozono y 24,61\% en $\Delta$. Las variantes de inicializaci\'on at\'omica tambi\'en generan sensibilidades importantes.

El resultado es un marco mecan\'istico validado num\'ericamente en un horizonte finito, no una climatolog\'ia validada mediante observaciones. La inicializaci\'on aproximada, los perfiles prescritos y la ausencia de transporte limitan su interpretaci\'on. Se priorizan mejores condiciones iniciales, transporte y comparaciones observacionales compatibles.

\chapter*{Acknowledgements}\addcontentsline{toc}{chapter}{Acknowledgements}
\missing{Optional personal acknowledgements; confirm whether required by the institution.}
